\chapter{Guarantees the composition layer supplies}

Two properties that no single expert has to possess, because the way experts
are called and connected creates them: an exact symmetry, and passivity.

\section{T10: symmetry averaging}

\paragraph{In plain words.}
A frozen expert does not respect a symmetry.  Call it once for every element of
the symmetry group, undo each transformation, and average.  The averaged map
respects the symmetry exactly, whatever the expert is: nonlinear, learned,
discontinuous.  The proof uses only that multiplying every group element by a
fixed one runs through the whole group again, which is why a set of operations
that is not a group does not work.

\begin{definition}[The group average]\label{def:average}
  \lean{Atlas.average}
  \leanok
  Let $G$ be a finite group acting on a set $V$ and acting linearly on a vector
  space $W$.  For any map $E:V\to W$,
  \[
    \tilde E(u)=\frac1{|G|}\sum_{g\in G}g^{-1}\,E(g\,u).
  \]
\end{definition}

\begin{theorem}[T10, symmetry averaging]\label{thm:T10}
  \lean{Atlas.average_equivariant}
  \leanok
  \uses{def:average}
  For every map $E$, every $h\in G$ and every $u\in V$,
  $\tilde E(h\,u)=h\,\tilde E(u)$.
\end{theorem}

\begin{proof}
  \uses{def:average}
  \leanok
  $\tilde E(hu)=\frac1{|G|}\sum_gg^{-1}E(ghu)$.  Substitute $g'=gh$, which runs
  over $G$ as $g$ does; then $g^{-1}=h\,g'^{-1}$, and $h$ comes out of the sum
  because the action on $W$ is linear.
\end{proof}

\begin{proposition}[Averaging leaves a symmetric map alone]\label{prop:T10-idem}
  \lean{Atlas.average_of_equivariant}
  \leanok
  \uses{def:average}
  If $E(gu)=gE(u)$ for all $g,u$, and $|G|$ is invertible in the scalars, then
  $\tilde E=E$.
\end{proposition}

\begin{proof}
  \uses{def:average}
  \leanok
  Every term of the sum is $E(u)$.
\end{proof}

\begin{example}[A set that is not a group]\label{ex:T10-nongroup}
  \lean{Atlas.average_over_non_group_not_equivariant}
  \leanok
  Let $M_x(x,y)=(-x,y)$ and $M_y(x,y)=(x,-y)$ on $\R^2$.  The set
  $\{\mathrm{id},M_x,M_y\}$ is not a group: $M_xM_y$ is missing.  The same
  average over this set is not $M_x$-equivariant: for $E(x,y)=(x+y,0)$ and
  $u=(0,1)$, the average at $M_xu=u$ is $(-\tfrac13,0)$, and $M_x$ of the
  average at $u$ is $(\tfrac13,0)$.
\end{example}

\begin{proof}
  \leanok
  $E(u)=(1,0)$, $M_xE(M_xu)=M_x(1,0)=(-1,0)$ and $M_yE(M_yu)=M_y(-1,0)=(-1,0)$,
  so the average at $u=M_xu$ is $\tfrac13(-1,0)$, and $M_x$ of it is
  $(\tfrac13,0)\ne(-\tfrac13,0)$.
\end{proof}

\section{T11: passive agents joined by the port rule}

\paragraph{In plain words.}
A port carries an effort and a flow whose inner product is power.  The port
algebra joins two ports by one rule: equal efforts, opposite flows.  Then the
power entering one port is exactly the power leaving its partner, so over all
connected ports the powers cancel, whatever the graph.  If every agent is
passive (its stored energy grows by no more than what entered through its
ports), the whole system is passive with respect to its open ports.  And if
every agent is \emph{incrementally} passive, one coupled step does not increase
the distance between two runs: the constant $L\le1$ of the master bound.

\begin{lemma}[The connection rule conserves power]\label{lem:T11-rule}
  \lean{Atlas.port_rule_power}
  \leanok
  In a real inner-product space, if $e_B=e_A$ and $f_B=-f_A$ then
  $\langle e_B,f_B\rangle=-\langle e_A,f_A\rangle$.
\end{lemma}

\begin{proof}
  \leanok
  Linearity of the inner product in its second argument.
\end{proof}

\begin{lemma}[The junction]\label{lem:T11-junction}
  \lean{Atlas.junction_power_eq_zero}
  \leanok
  Let $P$ be a finite set of ports, $\pi:P\to P$ a bijection pairing each port
  with its partner, and $p:P\to\R$ with $p(\pi(x))=-p(x)$ for all $x$.  Then
  $\sum_{x\in P}p(x)=0$.
\end{lemma}

\begin{proof}
  \leanok
  $\sum_xp(x)=\sum_xp(\pi(x))=-\sum_xp(x)$.
\end{proof}

\begin{theorem}[T11, passivity composes]\label{thm:T11}
  \lean{Atlas.interconnection_passive}
  \leanok
  \uses{lem:T11-junction}
  Let each port belong to one of finitely many agents.  Agent $i$ stores $H_i$
  before a step and $H_i'$ after it, receives the power of its own connected
  ports and $\mathrm{ext}_i$ through its open ports, and is passive:
  $H_i'\le H_i+\sum_{x\text{ of }i}p(x)+\mathrm{ext}_i$.  If the connected ports
  obey $p(\pi(x))=-p(x)$, then
  \[
    \sum_iH_i'\le\sum_iH_i+\sum_i\mathrm{ext}_i .
  \]
\end{theorem}

\begin{proof}
  \uses{lem:T11-junction}
  \leanok
  Sum the agents' inequalities.  The port powers, summed over agents, are the
  sum over all ports, which is zero by Lemma~\ref{lem:T11-junction}.
\end{proof}

\begin{corollary}[T11, incremental passivity gives a non-expansive step]\label{cor:T11-nonexpansive}
  \lean{Atlas.interconnection_nonexpansive}
  \leanok
  \uses{thm:T11}
  Let $u_i,v_i$ be agent $i$'s states in two runs before a step and $u_i',v_i'$
  after it, and let $\Delta p(x)$ be the power pairing of the differences of
  the two runs' efforts and flows at port $x$.  If
  $\norm{u_i'-v_i'}^2\le\norm{u_i-v_i}^2+2\sum_{x\text{ of }i}\Delta p(x)$ for
  every agent and $\Delta p(\pi(x))=-\Delta p(x)$, then
  \[
    \sum_i\norm{u_i'-v_i'}^2\le\sum_i\norm{u_i-v_i}^2 .
  \]
\end{corollary}

\begin{proof}
  \uses{thm:T11}
  \leanok
  Theorem~\ref{thm:T11} with $H_i=\norm{u_i-v_i}^2$, the powers $2\Delta p$, and
  no external term.
\end{proof}
